\documentclass[9pt,a4paper]{extarticle}

\usepackage[utf8]{inputenc}
\usepackage[T5]{fontenc}
\usepackage[vietnamese]{babel}
\usepackage[a4paper,top=6mm,bottom=6mm,left=6mm,right=6mm,columnsep=4mm]{geometry}
\usepackage{multicol}
\usepackage{enumitem}
\usepackage{titlesec}
\usepackage{microtype}
\usepackage{xcolor}
\usepackage{lmodern}
\usepackage{parskip}
\usepackage{array}
\usepackage{tabularx}

\pagestyle{empty}
\setlength{\parindent}{0pt}
\setlength{\parskip}{1pt}
\setlength{\columnseprule}{0pt}

\linespread{0.92}

\titleformat{\section}{\bfseries\fontsize{9.2}{10}\selectfont}{}{0pt}{}
\titleformat{\subsection}{\bfseries\fontsize{8.6}{9.2}\selectfont}{}{0pt}{}
\titlespacing*{\section}{0pt}{2pt}{1pt}
\titlespacing*{\subsection}{0pt}{1pt}{0.5pt}

\setlist[itemize]{leftmargin=*,topsep=1pt,itemsep=0.5pt,parsep=0pt,partopsep=0pt}
\setlist[enumerate]{leftmargin=*,topsep=1pt,itemsep=0.5pt,parsep=0pt,partopsep=0pt}

\begin{document}
\small
\begin{multicols*}{2}

%% ============================================================
\section*{Câu 1: Trình bày vắn tắt sự ra đời và phát triển của triết học Mác -- Lênin. Tại sao nói sự ra đời của triết học Mác -- Lênin là một bước ngoặt cách mạng trong lịch sử triết học? Trong thời đại hiện nay, triết học Mác -- Lênin có vai trò như thế nào?}

Sự ra đời và phát triển của triết học Mác -- Lênin là một quá trình tất yếu lịch sử, đánh dấu sự thay đổi căn bản về chất trong lịch sử tư tưởng nhân loại. Dưới đây là nội dung trình bày vắn tắt theo yêu cầu của bạn:

\subsection*{1. Sự ra đời và phát triển của triết học Mác -- Lênin}

\textbf{Điều kiện và tiền đề ra đời (những năm 40 của thế kỷ XIX):}

\begin{itemize}
\item \textbf{Điều kiện kinh tế - xã hội:} Sự phát triển mạnh mẽ của phương thức sản xuất tư bản chủ nghĩa dẫn đến sự bần cùng hóa giai cấp vô sản và thúc đẩy cuộc đấu tranh của giai cấp này chuyển từ tự phát sang tự giác.
\item \textbf{Tiền đề khoa học tự nhiên:} Dựa trên ba phát minh vĩ đại là Học thuyết tế bào, Thuyết tiến hóa của Đác-uyn và Định luật bảo toàn và chuyển hóa năng lượng.
\item \textbf{Tiền đề lý luận:} Kế thừa trực tiếp phép biện chứng của Hê-ghen và chủ nghĩa duy vật của Phoi-ơ-bắc trong nền triết học cổ điển Đức.
\end{itemize}

\textbf{Các giai đoạn phát triển:}

\begin{itemize}
\item \textbf{Giai đoạn C. Mác và Ph. Ăng-ghen:} Từ việc chuyển biến thế giới quan (1842--1843) đến việc đề xuất các nguyên lý duy vật biện chứng và duy vật lịch sử (1844--1848), sau đó tiếp tục bổ sung và phát triển lý luận (sau 1848).
\item \textbf{Giai đoạn V.I. Lênin:} Bảo vệ và phát triển triết học Mác trong điều kiện thực tiễn mới (chủ nghĩa đế quốc và cách mạng vô sản), đưa triết học Mác lên một giai đoạn phát triển mới.
\item \textbf{Giai đoạn hiện nay:} Tiếp tục được các Đảng Cộng sản và công nhân bổ sung, phát triển dựa trên thực tiễn sinh động của thời đại.
\end{itemize}

\subsection*{2. Tại sao sự ra đời của triết học Mác -- Lênin là một bước ngoặt cách mạng?}

Sự ra đời này được coi là một cuộc cách mạng vĩ đại trong lịch sử triết học vì những lý do sau:

\begin{itemize}
\item \textbf{Thay đổi mục đích của triết học:} Chuyển từ việc chủ yếu là \textbf{``giải thích'' thế giới} sang không chỉ giải thích mà còn tập trung vào nhiệm vụ \textbf{``cải tạo'' thế giới} bằng hành động cách mạng.
\item \textbf{Thống nhất hữu cơ giữa chủ nghĩa duy vật và phép biện chứng:} Khắc phục sự tách rời trước đó để tạo ra \textbf{chủ nghĩa duy vật biện chứng}, làm cho chủ nghĩa duy vật trở nên hoàn bị và phép biện chứng trở thành khoa học.
\item \textbf{Sáng tạo ra chủ nghĩa duy vật lịch sử:} Đây là biểu hiện vĩ đại nhất, kết thúc thời kỳ nghiên cứu lịch sử xã hội một cách tùy tiện, mở ra khả năng nghiên cứu xã hội như một quá trình lịch sử -- tự nhiên.
\item \textbf{Xác lập mối quan hệ đúng đắn với các khoa học cụ thể:} Chấm dứt quan niệm triết học là ``khoa học của các khoa học'', xác lập vai trò của triết học là \textbf{thế giới quan và phương pháp luận} chung nhất cho khoa học.
\item \textbf{Trở thành vũ khí lý luận của giai cấp vô sản:} Lần đầu tiên trong lịch sử, giai cấp vô sản có một học thuyết khoa học soi đường cho cuộc đấu tranh tự giải phóng mình và giải phóng nhân loại.
\end{itemize}

\subsection*{3. Vai trò của triết học Mác -- Lênin trong thời đại hiện nay}

Trong bối cảnh thế giới hiện nay, triết học Mác -- Lênin giữ những vai trò then chốt:

\begin{itemize}
\item \textbf{Là thế giới quan khoa học và phương pháp luận hiệu quả:} Cung cấp ``công cụ nhận thức vĩ đại'' để nhân loại tiến bộ giải quyết các vấn đề phức tạp của thời đại như toàn cầu hóa, kinh tế tri thức và cách mạng khoa học công nghệ hiện đại.
\item \textbf{Cơ sở lý luận cho công cuộc xây dựng chủ nghĩa xã hội:} Là nền tảng tư tưởng, kim chỉ nam cho các Đảng Cộng sản (như Đảng Cộng sản Việt Nam) trong việc hoạch định đường lối đổi mới và phát triển đất nước.
\item \textbf{Định hướng giá trị nhân văn:} Học thuyết về giải phóng con người vẫn còn nguyên giá trị, là nguồn cổ vũ cho cuộc đấu tranh vì tự do, công bằng và sự phát triển toàn diện của mỗi cá nhân.
\item \textbf{Hệ thống lý luận mở:} Triết học Mác -- Lênin không phải là giáo điều mà luôn đòi hỏi được \textbf{bổ sung, hoàn thiện và phát triển} từ thực tiễn sinh động, giúp nhân loại tiếp tục tìm kiếm và nhận thức chân lý.
\end{itemize}

%% ============================================================
\section*{Câu 2: Phân tích cơ sở lý luận và yêu cầu phương pháp luận của nguyên tắc phát triển. Sự vận dụng của C.Mác nguyên tắc này vào nghiên cứu quá trình vận động, phát triển xã hội loài người. Khái quát sự ra đời và phát triển của lĩnh vực trí tuệ nhân tạo.}

\subsection*{1. Nguyên tắc phát triển: Cơ sở lý luận và yêu cầu phương pháp luận}

\textbf{Cơ sở lý luận:}

\begin{itemize}
\item \textbf{Định nghĩa:} Phát triển là một phạm trù triết học dùng để chỉ quá trình vận động của sự vật theo khuynh hướng đi lên: từ trình độ thấp đến trình độ cao, từ đơn giản đến phức tạp, từ kém hoàn thiện đến hoàn thiện hơn.
\item \textbf{Tính chất:} Phát triển mang tính \textbf{khách quan} (là nguồn gốc nằm trong chính bản thân sự vật), tính \textbf{phổ biến} (diễn ra ở mọi lĩnh vực tự nhiên, xã hội và tư duy) và tính \textbf{đa dạng, phong phú}.
\item \textbf{Cơ chế:} Phát triển không diễn ra theo đường thẳng mà theo \textbf{đường xoắn ốc}, bao hàm cả những bước thụt lùi tạm thời. Sự phát triển là kết quả của quá trình tích lũy dần dần về lượng dẫn đến bước nhảy về chất, thông qua việc giải quyết các mâu thuẫn nội tại và trải qua các lần phủ định biện chứng.
\end{itemize}

\textbf{Yêu cầu phương pháp luận (Nguyên tắc phát triển):}

\begin{itemize}
\item \textbf{Khi nhận thức:} Phải đặt sự vật trong \textbf{sự vận động}, phát hiện các giai đoạn tồn tại và xu hướng biến đổi của nó trong tương lai.
\item \textbf{Trong hoạt động thực tiễn:} Phải nhạy bén với cái mới, ủng hộ cái mới phù hợp với quy luật dù nó còn non yếu. Phải biết thực hiện các bước nhảy cục bộ để chuẩn bị cho bước nhảy toàn bộ.
\item \textbf{Phê phán:} Cần khắc phục tư duy siêu hình, thái độ bảo thủ, trì trệ, định kiến hoặc giáo điều trong nhận thức và hành động.
\end{itemize}

\subsection*{2. Sự vận dụng nguyên tắc phát triển của C. Mác vào nghiên cứu xã hội}

C. Mác đã vận dụng nguyên tắc phát triển để xây dựng \textbf{quan niệm duy vật về lịch sử}, tạo ra bước ngoặt trong nghiên cứu xã hội loài người:

\begin{itemize}
\item \textbf{Tiến trình lịch sử -- tự nhiên:} C. Mác khẳng định sự phát triển của các hình thái kinh tế -- xã hội là một \textbf{quá trình lịch sử -- tự nhiên}. Xã hội không phải là một khối kết tinh vững chắc mà là một cơ thể luôn biến đổi và vận động theo các quy luật khách quan.
\item \textbf{Động lực phát triển:} Nguồn gốc sâu xa của sự phát triển xã hội là sự phát triển của \textbf{lực lượng sản xuất}. Khi lực lượng sản xuất phát triển đến một trình độ nhất định sẽ mâu thuẫn với quan hệ sản xuất lỗi thời, đòi hỏi một cuộc cách mạng kinh tế để thay thế quan hệ sản xuất cũ bằng quan hệ sản xuất mới phù hợp hơn, từ đó thúc đẩy toàn bộ kiến trúc thượng tầng thay đổi.
\item \textbf{Sự thay thế các hình thái kinh tế -- xã hội:} Lịch sử loài người là lịch sử thay thế tuần tự hoặc nhảy vọt của các phương thức sản xuất, từ Cộng sản nguyên thủy, Chiếm hữu nô lệ, Phong kiến, Tư bản chủ nghĩa đến Cộng sản chủ nghĩa.
\item \textbf{Giải phóng con người:} C. Mác coi mục đích tự thân của sự phát triển xã hội là sự \textbf{giải phóng và phát triển con người toàn diện}.
\end{itemize}

\subsection*{3. Khái quát sự ra đời và phát triển của lĩnh vực Trí tuệ nhân tạo (AI)}

Dựa trên bối cảnh cách mạng khoa học và công nghệ được mô tả trong tài liệu:

\begin{itemize}
\item \textbf{Bối cảnh ra đời:} AI xuất hiện và phát triển mạnh mẽ trong làn sóng của \textbf{cuộc cách mạng khoa học -- công nghệ hiện đại} (từ những năm 80 của thế kỷ XX đến nay). Đây là giai đoạn khoa học trở thành lực lượng sản xuất trực tiếp, dẫn dắt công nghệ và sản xuất.
\item \textbf{Bản chất phát triển:} AI là sản phẩm của \textbf{trí tuệ con người}, được tạo ra nhằm thay thế và hỗ trợ các chức năng trí tuệ của con người trong quá trình sản xuất và đời sống. Nó dựa trên các thành tựu của toán học, logic học, tin học và vật lý học hiện đại.
\item \textbf{Vị thế hiện nay:} AI được xác định là một \textbf{công nghệ chiến lược}, công nghệ cốt lõi trong cuộc Cách mạng công nghiệp lần thứ tư. Nó tham gia vào việc tối ưu hóa quản trị quốc gia, phát triển kinh tế số và thay đổi hoàn toàn phương thức tương tác của con người với thế giới.
\item \textbf{Định hướng tại Việt Nam:} Nhà nước Việt Nam đặt mục tiêu đến năm 2030 sẽ nằm trong \textbf{nhóm 3 nước dẫn đầu khu vực Đông Nam Á} về nghiên cứu và phát triển AI. AI được ứng dụng dựa trên dữ liệu lớn (Big Data) để trở thành tư liệu sản xuất chính trong kỷ nguyên mới.
\end{itemize}

%% ============================================================
\section*{Câu 3: Chỉ ra sự khác biệt cơ bản của mâu thuẫn lôgích và mâu thuẫn biện chứng. Giải thích luận điểm: \textit{Xét về thực chất, phép biện chứng là học thuyết về mâu thuẫn biện chứng}. Mâu thuẫn biện chứng nào đang chi phối sự vận động và phát triển lĩnh vực trí tuệ nhân tạo?}

\subsection*{1. Sự khác biệt cơ bản giữa mâu thuẫn lôgích và mâu thuẫn biện chứng}

Dựa trên các tài liệu trình chiếu, sự khác biệt giữa hai loại mâu thuẫn này được thể hiện qua ba phương diện chính:

\begin{center}
\renewcommand{\arraystretch}{1.2}
\begin{tabular}{|p{1.5cm}|p{3cm}|p{3cm}|}
\hline
\textbf{Đặc điểm} & \textbf{Mâu thuẫn biện chứng} & \textbf{Mâu thuẫn lôgích} \\
\hline
\textbf{Biểu hiện} & Là sự \textbf{thống nhất và đấu tranh của các mặt đối lập} vốn có trong bản thân sự vật. & Là sự kết hợp, khẳng định hoặc phủ định hai tư tưởng trái ngược nhau trong cùng một quá trình tư duy. \\
\hline
\textbf{Tính chất} & Mang tính \textbf{khách quan, phổ biến và đa dạng}, tồn tại độc lập với ý chí con người. & Mang tính \textbf{chủ quan}, không phổ biến, thường xuất hiện do sai sót trong lập luận. \\
\hline
\textbf{Vai trò} & Là \textbf{nguồn gốc và động lực} của sự vận động, phát triển của vạn vật trong thế giới. & Đưa nhận thức đến sai lầm và khiến tư duy rơi vào tình trạng bế tắc. \\
\hline
\end{tabular}
\end{center}

\subsection*{2. Giải thích luận điểm: ``Xét về thực chất, phép biện chứng là học thuyết về mâu thuẫn biện chứng''}

Luận điểm này (thường được gắn với V.I. Lênin) khẳng định rằng hạt nhân của phép biện chứng chính là quy luật mâu thuẫn vì:

\begin{itemize}
\item \textbf{Mâu thuẫn là nguồn gốc của sự phát triển:} Tài liệu khẳng định mọi sự vật, hiện tượng đều chứa đựng các mặt đối lập; sự đấu tranh giữa chúng là động lực thúc đẩy sự vật cũ mất đi, sự vật mới ra đời.
\item \textbf{Chi phối các quy luật khác:} Quy luật mâu thuẫn giải thích cho \textit{tại sao} có sự vận động. Các quy luật khác như Lượng -- Chất (cách thức vận động) và Phủ định của phủ định (khuynh hướng vận động) đều dựa trên việc giải quyết các mâu thuẫn nội tại để thực hiện bước nhảy về chất hoặc phủ định biện chứng.
\item \textbf{Tính phổ biến:} Mâu thuẫn biện chứng tồn tại trong mọi lĩnh vực từ tự nhiên, xã hội đến tư duy, là đối tượng nghiên cứu trung tâm để hiểu về sự biến đổi của thế giới.
\end{itemize}

\subsection*{3. Mâu thuẫn biện chứng chi phối sự vận động và phát triển của lĩnh vực Trí tuệ nhân tạo (AI)}

Trong bối cảnh cách mạng công nghệ hiện đại, lĩnh vực AI đang bị chi phối bởi các mâu thuẫn biện chứng sâu sắc sau:

\begin{itemize}
\item \textbf{Mâu thuẫn giữa Trí tuệ con người (nguyên bản) và Trí tuệ nhân tạo (mô phỏng):} AI là sản phẩm từ trí tuệ con người nhưng lại đang tiến tới hỗ trợ và thay thế các chức năng trí tuệ của con người. Mâu thuẫn này thúc đẩy con người không ngừng sáng tạo ra các thế hệ AI mới mạnh mẽ hơn, đồng thời đòi hỏi con người phải nâng cao năng lực bản thân để không bị ``lấn át'' bởi chủ nghĩa kỹ trị.
\item \textbf{Mâu thuẫn giữa Lực lượng sản xuất (AI) và Quan hệ sản xuất (Thể chế/Pháp luật):} AI phát triển như một lực lượng sản xuất hiện đại và thần tốc, nhưng hệ thống thể chế, pháp luật và quan niệm xã hội thường biến đổi chậm hơn (tính lạc hậu tương đối của ý thức xã hội). Sự xung đột này đòi hỏi phải ``đổi mới tư duy xây dựng pháp luật'', loại bỏ tư duy ``không quản được thì cấm'' để tạo điều kiện cho AI phát triển phù hợp với quản trị quốc gia.
\item \textbf{Mâu thuẫn giữa Lợi ích kinh tế và An ninh/Đạo đức:} Sự phát triển AI mang lại giá trị kinh tế khổng lồ nhưng cũng tạo ra những thách thức về an ninh dữ liệu, an toàn thông tin và bảo vệ quyền riêng tư của con người. Việc giải quyết mâu thuẫn này là động lực để xây dựng các bộ quy tắc ứng xử trên không gian mạng và phát triển AI theo hướng nhân văn.
\end{itemize}

%% ============================================================
\section*{Câu 4: Cái mới? Cái cũ? Chứng minh bằng tư duy lý luận, minh họa bằng tài liệu thực tiễn (trong lĩnh vực công nghệ thông tin...) luận điểm: ``Cuộc đấu tranh giữa cái mới và cái cũ là một quá trình khó khăn, phức tạp và lâu dài; Cái mới thường thất bại tạm thời nhưng cuối cùng nó sẽ chiến thắng cái cũ''. Từ đó rút ra những kết luận về mặt phương pháp luận dùng trong lĩnh vực hoạt động công nghệ thông tin của Anh/Chị.}

\subsection*{1. Khái niệm Cái cũ và Cái mới}

\begin{itemize}
\item \textbf{Cái cũ:} Là những gì đã và đang tồn tại nhưng \textbf{không còn phù hợp với quy luật khách quan} hoặc điều kiện thực tế của thời đại. Cái cũ thường chứa đựng các yếu tố thoái bộ, tiêu cực, sức sống giảm dần và tất yếu sẽ bị phủ định.
\item \textbf{Cái mới:} Là những gì chưa từng tồn tại trước đó, \textbf{phù hợp với quy luật và xu thế của thời đại}. Cái mới chứa đựng các yếu tố tiến bộ, tích cực, sức sống ngày càng lớn mạnh và cuối cùng sẽ được khẳng định.
\end{itemize}

\subsection*{2. Chứng minh tư duy lý luận về cuộc đấu tranh giữa cái mới và cái cũ}

Luận điểm: \textit{``Cuộc đấu tranh giữa cái mới và cái cũ là một quá trình khó khăn, phức tạp và lâu dài; Cái mới thường thất bại tạm thời nhưng cuối cùng nó sẽ chiến thắng cái cũ''} được lý giải qua các khía cạnh sau:

\begin{itemize}
\item \textbf{Tính khó khăn và phức tạp:} Phát triển không diễn ra theo đường thẳng mà theo \textbf{đường xoắn ốc}, bao hàm cả những bước thụt lùi tạm thời. Cái cũ thường gắn liền với những hệ thống, thiết chế và thói quen đã định hình lâu đời, nên nó có sức kháng cự mạnh mẽ.
\item \textbf{Sự thất bại tạm thời của cái mới:} Khi mới ra đời, cái mới thường còn non yếu, chưa hoàn thiện về cấu trúc và thiếu các điều kiện cần thiết để bộc lộ hết ưu thế. Ngược lại, cái cũ tuy đã lỗi thời nhưng vẫn còn giữ được những lực lượng nhất định che chở.
\item \textbf{Sự chiến thắng tất yếu của cái mới:} Vì cái mới phù hợp với quy luật khách quan, nó có \textbf{sức sống nội tại} mãnh liệt hơn. Thông qua quá trình phủ định biện chứng, cái mới kế thừa những hạt nhân tích cực của cái cũ đồng thời loại bỏ những gì cản trở, từ đó tạo ra bước nhảy về chất để khẳng định vị thế.
\end{itemize}

\subsection*{3. Minh họa bằng thực tiễn trong lĩnh vực Công nghệ thông tin}

\begin{itemize}
\item \textbf{Chuyển đổi số (Digital Transformation):} Đây là một ``cuộc cách mạng sâu sắc'' nhằm thay thế các quy trình quản lý thủ công, rời rạc (cái cũ) bằng hệ thống dữ liệu tập trung và tự động (cái mới). Ban đầu, các giải pháp số thường gặp sự phản ứng từ thói quen làm việc cũ hoặc lo ngại về an ninh dữ liệu. Tuy nhiên, với ưu thế vượt trội về hiệu suất, chuyển đổi số đang dần trở thành xu thế không thể đảo ngược.
\item \textbf{Trí tuệ nhân tạo (AI):} AI ra đời đối mặt với những định kiến và rào cản pháp lý ``không quản được thì cấm'' (cái cũ). Sự phát triển của AI có thể trải qua những giai đoạn thoái bộ khi gây ra các vấn đề về đạo đức hoặc an ninh mạng. Nhưng vì AI là công nghệ chiến lược phù hợp với lực lượng sản xuất hiện đại, nó tất yếu sẽ chiến thắng và trở thành tư liệu sản xuất chính.
\item \textbf{Mô hình Agile/DevOps:} Trong phát triển phần mềm, việc chuyển từ mô hình thác nước (Waterfall -- cũ) sang Agile (mới) ban đầu gặp nhiều thất bại do sự bảo thủ trong tư duy quản lý. Nhưng cuối cùng, Agile đã thắng thế vì nó đáp ứng được sự thay đổi nhanh chóng của thị trường phần mềm hiện đại.
\end{itemize}

\subsection*{4. Kết luận về mặt phương pháp luận cho hoạt động CNTT}

Từ lý luận trên, các chuyên gia CNTT cần rút ra những nguyên tắc phương pháp luận sau:

\begin{enumerate}
\item \textbf{Nhạy bén và ủng hộ cái mới:} Phải có tư duy cởi mở, vận dụng sáng tạo và sẵn sàng thí điểm những công nghệ mới (như AI, IoT, Cloud) dù ban đầu còn nhiều rủi ro.
\item \textbf{Dũng cảm phủ định cái cũ lỗi thời:} Kiên quyết từ bỏ những công nghệ lạc hậu, quy trình trì trệ hoặc tư duy quản lý kiểu ``bao cấp thông tin'' để tạo không gian cho đổi mới sáng tạo.
\item \textbf{Kiên trì và không bi quan trước thất bại:} Hiểu rằng các dự án công nghệ mới có thể gặp trục trặc tạm thời (độ trễ, lỗi hệ thống). Cần chấp nhận rủi ro và đầu tư mạo hiểm để đạt được bước nhảy vọt trong tương lai.
\item \textbf{Kế thừa và phát triển:} Khi xây dựng hệ thống mới, không nên ``phủ định sạch trơn'' mà cần kế thừa những dữ liệu, kinh nghiệm quý báu từ hệ thống cũ, đồng thời làm giàu chúng bằng công nghệ hiện đại.
\item \textbf{Gắn lý luận với thực tiễn:} Mọi giải pháp CNTT phải xuất phát từ nhu cầu thực tế của người dân và doanh nghiệp, lấy hiệu quả ứng dụng làm thước đo cao nhất cho chân lý.
\end{enumerate}

%% ============================================================
\section*{Câu 5: Chứng minh luận điểm của V.I.Lênin: ``Ý thức con người không chỉ phản ánh thế giới mà còn góp phần sáng tạo ra thế giới''. Tại sao nói, trong thời đại chúng ta, ý thức có vai trò rất to lớn, thậm chí quyết định, trong nhiều lĩnh vực hoạt động của con người? Quan điểm của Anh/Chị về vai trò và sự tác động của trí tuệ nhân tạo trong đời sống xã hội?}

Luận điểm của V.I.Lênin: \textbf{``Ý thức con người không chỉ phản ánh thế giới mà còn góp phần sáng tạo ra thế giới''} là một khẳng định quan trọng về tính năng động, sáng tạo của ý thức trong mối quan hệ với vật chất. Dưới đây là phần chứng minh luận điểm này, lý giải vai trò của ý thức trong thời đại mới và quan điểm về trí tuệ nhân tạo (AI):

\subsection*{1. Chứng minh luận điểm của V.I.Lênin}

Ý thức không phải là sự phản ánh thụ động, máy móc mà là sự \textbf{phản ánh năng động, sáng tạo} hiện thực khách quan vào bộ óc con người.

\begin{itemize}
\item \textbf{Tính sáng tạo trong phản ánh:} Ý thức là hình ảnh chủ quan của thế giới khách quan, trong đó vật chất được ``di chuyển'' vào đầu óc con người và được cải biến đi ở trong đó. Con người không chỉ chụp lại cái có sẵn mà còn biết \textbf{mô hình hóa đối tượng} trong tư duy dưới dạng các hình tượng, tư tưởng để tạo ra một ``thế giới tinh thần'' phong phú.
\item \textbf{Sáng tạo ra thế giới thông qua thực tiễn:} Ý thức tự thân không làm biến đổi thế giới, nhưng nó tham gia \textbf{chỉ đạo hành vi} con người trong hoạt động thực tiễn. Thông qua thực tiễn, những mô hình tư tưởng được \textbf{vật chất hóa} để tạo ra ``giới tự nhiên thứ hai'' -- đó là xã hội và các thành tựu văn minh mà nếu không có ý thức con người thì chúng không thể tồn tại.
\item \textbf{Sức mạnh của lý luận:} Khi các tư tưởng, lý luận khoa học (một dạng cao của ý thức) thâm nhập vào quần chúng, nó sẽ biến thành \textbf{lực lượng vật chất} để cải tạo xã hội và tự nhiên.
\end{itemize}

\subsection*{2. Vai trò to lớn, quyết định của ý thức trong thời đại ngày nay}

Trong thời đại hiện nay, ý thức đóng vai trò then chốt, thậm chí quyết định sự thành bại trong nhiều lĩnh vực vì:

\begin{itemize}
\item \textbf{Khoa học trở thành lực lượng sản xuất trực tiếp:} Các sản phẩm của tư duy như tri thức, công nghệ số và đổi mới sáng tạo đã trở thành động lực chính cho sự phát triển kinh tế. Tri thức hiện nay là yếu tố cốt lõi quy định năng suất và hiệu quả của nền sản xuất hiện đại.
\item \textbf{Định hướng và quản trị quốc gia:} Ý thức dưới dạng \textbf{chiến lược, thể chế và chính sách} đóng vai trò dẫn dắt, tạo điều kiện cho sự bứt phá của quốc gia. Một tư duy pháp luật đổi mới, loại bỏ tư duy ``không quản được thì cấm'' là tiền đề tiên quyết để khuyến khích sáng tạo.
\item \textbf{Giải quyết các vấn đề toàn cầu:} Những thách thức như biến đổi khí hậu, an ninh mạng hay dịch bệnh đòi hỏi con người phải có \textbf{tư duy lý luận sắc bén} và tầm nhìn xa để hoạch định các chiến lược ứng xử nhân văn với thế giới xung quanh.
\item \textbf{Nguồn lực con người là vốn quý nhất:} Trong nền kinh tế tri thức, trình độ trí tuệ và phẩm chất đạo đức của người lao động mới là yếu tố quyết định giá trị thặng dư, chứ không chỉ là sức lao động cơ bắp đơn thuần.
\end{itemize}

\subsection*{3. Vai trò và sự tác động của trí tuệ nhân tạo (AI) trong đời sống xã hội}

Từ góc độ triết học và thực tiễn phát triển công nghệ hiện nay, vai trò của AI có thể được nhìn nhận như sau:

\begin{itemize}
\item \textbf{Sản phẩm vĩ đại của trí tuệ con người:} AI là chứng minh cho sức mạnh sáng tạo vô hạn của ý thức con người; nó không tự nhiên có mà là kết quả của quá trình \textbf{sản xuất tri thức} và sáng tạo kỹ thuật.
\item \textbf{Công cụ giải phóng và nâng cao năng lực:} AI được xác định là một \textbf{công nghệ chiến lược}, giúp tối ưu hóa quản trị, thúc đẩy kinh tế số và thay thế con người trong các quy trình phức tạp, từ đó giải phóng con người khỏi những nhịp điệu làm việc máy móc.
\item \textbf{Động lực cho sự phát triển bền vững:} AI hỗ trợ nghiên cứu cơ bản, dự báo thiên tai và bảo vệ môi trường, giúp con người tiến gần hơn đến mục tiêu xây dựng một xã hội nhân văn và hiện đại.
\item \textbf{Thách thức và giới hạn:} Tuy AI có khả năng xử lý thông tin vượt trội, nhưng nó vẫn là công cụ hỗ trợ và \textbf{không thể thay thế hoàn toàn bản tính người} hay sức mạnh sáng tạo gốc của chủ thể có ý thức. Sự phát triển AI cần được định hướng bởi một thế giới quan đúng đắn để tránh rơi vào ``chủ nghĩa kỹ trị'' thống trị xã hội.
\item \textbf{Định hướng chiến lược tại Việt Nam:} Nhà nước coi phát triển AI là nhiệm vụ đột phá, đặt mục tiêu Việt Nam nằm trong \textbf{nhóm 3 nước dẫn đầu Đông Nam Á} về nghiên cứu và phát triển trí tuệ nhân tạo vào năm 2030. Việc làm chủ AI dựa trên dữ liệu lớn sẽ giúp Việt Nam tự chủ về công nghệ và nâng cao năng lực cạnh tranh quốc tế.
\end{itemize}

%% ============================================================
\section*{Câu 6: Nhận thức? Chân lý? Theo quan điểm của Anh/Chị, nhận thức nhân loại là một quá trình phát hiện ngày càng nhiều sai lầm hay là tích lũy ngày càng nhiều chân lý? Tại sao? Lấy ví dụ minh họa.}

\subsection*{1. Khái niệm Nhận thức và Chân lý}

\begin{itemize}
\item \textbf{Nhận thức:} Là quá trình phản ánh năng động, sáng tạo thế giới khách quan vào bộ óc con người; là quá trình xâm nhập sâu rộng của lý trí con người vào thế giới xung quanh để tìm hiểu, nắm bắt các cấp độ quy luật và bản chất của đối tượng. Nhận thức diễn ra theo con đường biện chứng: từ trực quan sinh động đến tư duy trừu tượng, và từ tư duy trừu tượng đến thực tiễn.
\item \textbf{Chân lý:} Là tri thức có nội dung phù hợp với khách thể mà nó phản ánh, đồng thời đã được thực tiễn kiểm nghiệm. Chân lý có tính khách quan, tính cụ thể và là một quá trình vận động từ chân lý tương đối (phản ánh đúng nhưng chưa đầy đủ) tiến dần đến chân lý tuyệt đối (phản ánh đầy đủ, hoàn chỉnh).
\end{itemize}

\subsection*{2. Nhận thức nhân loại: Phát hiện sai lầm hay tích lũy chân lý?}

Theo quan điểm triết học được trình bày trong các nguồn tài liệu, nhận thức nhân loại \textbf{vừa là quá trình tích lũy chân lý, vừa là quá trình phát hiện và vượt qua sai lầm}. Hai mặt này không tách rời mà thống nhất biện chứng với nhau:

\begin{itemize}
\item \textbf{Tích lũy chân lý (Góc nhìn Mác-xít):} Nhận thức không đi đến chân lý tuyệt đỉnh ngay lập tức mà vận động xuyên qua các ``chân lý tương đối''. Mỗi giai đoạn lịch sử bồi đắp thêm những hiểu biết mới, làm cho bức tranh về thế giới ngày càng chính xác và đầy đủ hơn. Các giá trị tri thức bền vững được kế thừa và phát triển qua các thời đại.
\item \textbf{Phát hiện sai lầm (Góc nhìn của K. Popper):} Theo chủ nghĩa phủ chứng của Karl Popper, khoa học không truy tìm tính chân lý để xác chứng mà là \textbf{truy tìm sai lầm để phủ chứng} lý luận. Ông cho rằng tri thức nhân loại luôn không đầy đủ, và khoa học tiến bộ nhanh hơn khi các lý luận sai lầm bị bác bỏ nhanh chóng để nhường chỗ cho những giả thuyết mới tốt hơn.
\item \textbf{Thực chất của quá trình:} Nhận thức là một \textbf{đường xoắn ốc}. Cái mới ra đời không phải là sự phủ định sạch trơn cái cũ, mà là phát hiện ra những giới hạn (sai lầm do tính hạn chế của thời đại) của cái cũ để xác lập một hệ thống tri thức hoàn thiện hơn. Vì vậy, việc phát hiện sai lầm chính là ``phương tiện'' để nhân loại tích lũy thêm những ``hạt nhân'' chân lý sâu sắc hơn.
\end{itemize}

\subsection*{3. Ví dụ minh họa}

Sự phát triển của \textbf{Vật lý học} là minh chứng điển hình cho quá trình này:

\begin{itemize}
\item \textbf{Bức tranh cơ học của Newton (Cái cũ/Chân lý tương đối):} Trong thời gian dài, các quy luật của Newton về không gian và thời gian tuyệt đối được coi là chân lý vĩnh cửu.
\item \textbf{Phát hiện sai lầm/hạn chế:} Đến cuối thế kỷ XIX, các phát minh về điện tử và phóng xạ cho thấy cơ học Newton không còn giải thích được các hiện tượng ở tốc độ cận ánh sáng hoặc trong thế giới vi mô (cuộc khủng hoảng vật lý).
\item \textbf{Tích lũy chân lý mới (Thuyết tương đối Einstein):} Einstein đã chỉ ra những ``sai lầm'' (hạn chế) của Newton về không gian và thời gian để xây dựng nên một lý thuyết tổng quát hơn. Thuyết tương đối không xóa bỏ hoàn toàn cơ học Newton mà giữ lại nó như một trường hợp đặc biệt trong điều kiện tốc độ thấp, đồng thời tích lũy thêm những hiểu biết vạch thời đại về mối quan hệ giữa vật chất và không gian -- thời gian.
\end{itemize}

\textbf{Kết luận:} Nhận thức là hành trình nhân loại \textbf{dũng cảm phạm sai lầm và học tập từ chính những sai lầm đó} để không ngừng tiến gần hơn đến chân lý khách quan. Như V.I. Lênin đã nhận định: ``Chỉ có ai không làm gì mới không bao giờ sai lầm''.

%% ============================================================
\section*{Câu 7: Chỉ ra nội dung triết học trong câu nói của V.I.Lênin: \textit{``Vấn đề tìm hiểu xem tư duy của con người có thể đạt tới chân lý khách quan hay không, hoàn toàn không phải là một vấn đề lý luận mà là một vấn đề thực tiễn. Chính trong thực tiễn mà con người phải chứng minh chân lý''.}}

Câu nói này (được V.I. Lênin trích dẫn lại từ tư tưởng của C. Mác trong tác phẩm \textit{Luận cương về Phoi-ơ-bắc}) được coi là \textbf{tuyên ngôn của một nền triết học hành động và triết học thực tiễn}. Nội dung triết học sâu sắc của câu nói tập trung vào các điểm chính sau đây:

\subsection*{1. Thực tiễn là tiêu chuẩn khách quan và tối cao của chân lý}

Triết học Mác -- Lênin khẳng định rằng việc kiểm tra tính đúng đắn của tư duy không thể thực hiện bằng các lập luận lý thuyết thuần túy hay tranh luận kinh viện.

\begin{itemize}
\item \textbf{Chân lý} là tri thức có nội dung phù hợp với hiện thực khách quan. Để biết một ý tưởng có đạt tới ``chân lý khách quan'' hay không, con người không thể chỉ dừng lại ở suy tư mà phải đưa nó vào \textbf{hoạt động thực tiễn}.
\item \textbf{Thực tiễn chính là ``thước đo''}: Chỉ thông qua thực tiễn, những mô hình tư tưởng mới được ``vật chất hóa'' để chứng minh tính hiện thực, sức mạnh và ``tính trần tục'' (tính thực tế) của tư duy. Nếu lý luận dẫn đến kết quả thành công trong thực tiễn, nó chứa đựng chân lý; nếu thất bại, nó cần được điều chỉnh.
\end{itemize}

\subsection*{2. Triết học cải tạo thế giới bằng hành động}

Nội dung câu nói xác lập một bước ngoặt về \textbf{mục đích của triết học}.

\begin{itemize}
\item Trước đây, các nhà triết học chủ yếu tập trung vào việc \textbf{``giải thích'' thế giới} bằng nhiều cách khác nhau.
\item C. Mác và V.I. Lênin khẳng định: Vấn đề quan trọng hơn là \textbf{``cải tạo'' thế giới} bằng cách mạng. Triết học phải lấy sinh khí từ thực tiễn và quay lại phục vụ thực tiễn để biến đổi hiện thực.
\end{itemize}

\subsection*{3. Sự thống nhất biện chứng giữa lý luận và thực tiễn}

Câu nói nhấn mạnh rằng nhận thức chân lý là một \textbf{vấn đề thực tiễn}, xác lập mối quan hệ không thể tách rời giữa lý luận và hành động.

\begin{itemize}
\item \textbf{Lý luận không được xa rời thực tiễn:} Lý luận mà không liên hệ với thực tiễn chỉ là lý luận suông, giáo điều và không có sức sống.
\item \textbf{Lý luận trở thành lực lượng vật chất:} Tư tưởng chỉ thực sự có sức mạnh khi nó thâm nhập vào quần chúng và được hiện thực hóa thông qua hoạt động thực tiễn của họ.
\item Thực tiễn là cơ sở, động lực và mục đích của nhận thức, đồng thời lý luận có vai trò soi đường, dẫn dắt thực tiễn để tránh sự ``mù quáng''.
\end{itemize}

\subsection*{4. Phê phán tính chất tư biện, trực quan của triết học cũ}

Câu nói này nhằm khắc phục khiếm khuyết lớn nhất của chủ nghĩa duy vật trước Mác (kể cả Phoi-ơ-bắc).

\begin{itemize}
\item Triết học cũ coi hiện thực hay cái cảm giác được chỉ là đối tượng để quan sát (trực quan) chứ không hiểu đó là kết quả của \textbf{hoạt động cảm giác của con người}, tức là thực tiễn.
\item Mác và Lênin khẳng định con người không chỉ phản ánh thế giới một cách thụ động mà còn là \textbf{chủ thể sáng tạo}, thông qua thực tiễn để làm biến đổi thế giới và chứng minh chân lý của mình.
\end{itemize}

\textbf{Tóm lại}, nội dung cốt lõi là xác lập \textbf{quan điểm thực tiễn} là quan điểm cơ bản, hàng đầu của lý luận nhận thức, khẳng định chân lý không phải là thứ để chiêm nghiệm mà là thứ phải được \textbf{con người chứng minh qua hành động thực tế}.

%% ============================================================
\section*{Câu 8: Chỉ ra nội dung triết học trong câu nói của C.Mác: \textit{``Vũ khí của sự phê phán cố nhiên không thể thay thế được sự phê phán của vũ khí, lực lượng vật chất chỉ có thể bị đánh đổ bằng lực lượng vật chất; nhưng lý luận cũng sẽ trở thành lực lượng vật chất, một khi nó thâm nhập vào quần chúng''.}}

Câu nói này của C. Mác (trích trong tác phẩm \textit{Góp phần phê phán triết học pháp quyền của Hêghen -- Lời mở đầu}) chứa đựng những nội dung triết học sâu sắc về mối quan hệ giữa lý luận và thực tiễn, cũng như vai trò của quần chúng nhân dân trong lịch sử:

\begin{itemize}
\item \textbf{Khẳng định vai trò quyết định của thực tiễn và lực lượng vật chất:} Mác chỉ rõ rằng ``vũ khí của sự phê phán'' (tức là những lý luận, tư tưởng) không thể thay thế được ``sự phê phán của vũ khí'' (tức là hành động thực tiễn, bạo lực cách mạng). Điều này nhấn mạnh rằng để thay đổi một hiện thực xã hội lỗi thời hay đánh đổ một hệ thống cai trị, con người không thể chỉ dùng lời nói hay tư tưởng mà phải sử dụng \textbf{lực lượng vật chất} tương ứng.
\item \textbf{Sức mạnh cải tạo thế giới của lý luận khoa học:} Tuy nhấn mạnh thực tiễn, nhưng Mác cũng không hạ thấp vai trò của lý luận. Ông khẳng định lý luận sẽ trở thành một \textbf{``lực lượng vật chất''} khi nó được quần chúng nhân dân tiếp thu và vận dụng. Khi những tư tưởng khoa học và cách mạng thâm nhập vào quần chúng, nó sẽ biến thành niềm tin, ý chí và sức mạnh dẫn dắt hành động thực tiễn của họ để cải tạo xã hội.
\item \textbf{Mối quan hệ biện chứng giữa triết học và giai cấp vô sản:} Nội dung câu nói này còn hàm ý về sự liên minh tất yếu: ``Giống như triết học thấy giai cấp vô sản là vũ khí vật chất của mình, giai cấp vô sản cũng thấy triết học là vũ khí tinh thần của mình''. Triết học (lý luận) cung cấp kim chỉ nam và mục tiêu, còn quần chúng nhân dân (lực lượng vật chất) là người thực thi lý tưởng đó vào đời sống.
\item \textbf{Bước ngoặt từ ``giải thích'' sang ``cải tạo'' thế giới:} Đây là tuyên ngôn cho một nền triết học hành động. Nó khắc phục khiếm khuyết của các trào lưu triết học cũ vốn chỉ dừng lại ở việc \textbf{giải thích thế giới} bằng các hệ thống tư tưởng trừu tượng, đồng thời khẳng định sứ mệnh của triết học Mác là phải tham gia trực tiếp vào công cuộc \textbf{cải tạo thế giới} bằng cách mạng thực tiễn.
\end{itemize}

\textbf{Tóm lại}, nội dung triết học cốt lõi là xác lập \textbf{quan điểm thực tiễn}, khẳng định sức mạnh to lớn của lý luận khi gắn liền với phong trào cách mạng của quần chúng, tạo ra sự thống nhất hữu cơ giữa sức mạnh trí tuệ và sức mạnh vật chất để biến đổi hiện thực.

%% ============================================================
\section*{Câu 9: Phân tích cơ sở lý luận và yêu cầu phương pháp luận của nguyên tắc thống nhất lý luận và thực tiễn. Sự vận dụng nguyên tắc này trong lĩnh vực công nghệ thông tin?}

\subsection*{1. Cơ sở lý luận của nguyên tắc thống nhất lý luận và thực tiễn}

Nguyên tắc này dựa trên \textbf{mối quan hệ biện chứng giữa thực tiễn và lý luận}, trong đó thực tiễn đóng vai trò quyết định và lý luận có tác động trở lại mạnh mẽ:

\begin{itemize}
\item \textbf{Vai trò quyết định của thực tiễn:} Thực tiễn là \textbf{cơ sở, động lực, mục đích} của nhận thức và là \textbf{tiêu chuẩn khách quan} duy nhất để kiểm tra chân lý. Lý luận không tự nhiên nảy sinh mà là sự tổng kết những kinh nghiệm của loài người, được hình thành và phát triển từ nhu cầu giải quyết các vấn đề do thực tiễn đặt ra.
\item \textbf{Vai trò của lý luận:} Lý luận là hệ thống tri thức mang tính khái quát, phổ quát về thế giới. Tuy thực tiễn mang tính thứ nhất, nhưng \textbf{lý luận có vai trò hướng dẫn, chỉ đạo thực tiễn}; nếu thực tiễn không có lý luận soi đường sẽ trở thành ``thực tiễn mù quáng''.
\item \textbf{Sự thống nhất biện chứng:} Lý luận và thực tiễn là hai mặt đối lập nhưng luôn thống nhất và bổ sung cho nhau. Lý luận chỉ có sức mạnh thực sự khi nó \textbf{thâm nhập vào quần chúng} và được vật chất hóa thông qua hoạt động thực tiễn để cải tạo thế giới. Ngược lại, lý luận phải liên tục thay đổi và phát triển cùng với sự biến đổi của thực tiễn để không trở nên lạc hậu, giáo điều.
\end{itemize}

\subsection*{2. Yêu cầu phương pháp luận}

Từ cơ sở lý luận trên, nguyên tắc này đặt ra các yêu cầu cụ thể sau:

\begin{itemize}
\item \textbf{Lý luận phải xuất phát từ thực tiễn:} Mọi quan điểm, học thuyết phải bám sát hiện thực, được tổng kết từ kinh nghiệm thực tiễn và phản ánh đúng các quy luật khách quan.
\item \textbf{Học đi đôi với hành:} Lý luận phải được áp dụng vào công việc thực tế, nếu chỉ học thuộc lòng sách vở mà không biết đem ra thực hành thì chỉ là ``lý luận suông'' và vô ích.
\item \textbf{Khắc phục các căn bệnh tư duy:}
  \begin{itemize}
  \item \textbf{Bệnh giáo điều:} Vận dụng lý luận một cách máy móc, rập khuôn mà không chú trọng đến điều kiện thực tiễn cụ thể.
  \item \textbf{Bệnh kinh nghiệm:} Tuyệt đối hóa kinh nghiệm thực tế, coi khinh lý luận, dẫn đến cách làm việc mò mẫm, không hiểu rõ toàn cục của sự phát triển.
  \end{itemize}
\item \textbf{Nói đi đôi với làm:} Đây là chuẩn mực đạo đức cách mạng quan trọng nhất, đòi hỏi lời nói phải đúng chủ trương, không hứa suông và phải được chứng minh bằng kết quả công việc cụ thể.
\end{itemize}

\subsection*{3. Sự vận dụng trong lĩnh vực Công nghệ thông tin (CNTT)}

Trong bối cảnh cách mạng công nghệ hiện nay, việc thống nhất lý luận và thực tiễn được thể hiện rõ nét qua các khía cạnh:

\begin{itemize}
\item \textbf{Khoa học dẫn dắt công nghệ:} Trong lĩnh vực CNTT, khoa học (lý luận) đóng vai trò dẫn đường và là động lực trực tiếp cho sự phát triển của công nghệ và sản xuất. Các lý thuyết về thuật toán, cấu trúc dữ liệu phải liên tục được kiểm nghiệm và tối ưu hóa thông qua các sản phẩm thực tế để đáp ứng nhu cầu xã hội.
\item \textbf{Chuyển đổi số quốc gia:} Đây được coi là một cuộc cách mạng sâu sắc, đòi hỏi \textbf{thể chế (lý luận/pháp luật) phải đi trước một bước} để dẫn dắt. Thực tiễn chuyển đổi số đòi hỏi phải đổi mới tư duy xây dựng pháp luật, loại bỏ tư duy ``không quản được thì cấm'' để tạo không gian cho đổi mới sáng tạo phát triển.
\item \textbf{Phát triển Trí tuệ nhân tạo (AI) và Dữ liệu:} Sự phát triển của AI không chỉ dựa trên các công thức toán học (lý luận) mà phải dựa trên \textbf{tiềm năng của dữ liệu lớn} (thực tiễn). Việc ứng dụng AI vào các ngành như y tế, giao thông, tài chính là quá trình đưa lý luận vào giải quyết các vấn đề cụ thể của đời sống, từ đó AI lại được hoàn thiện hơn thông qua việc học tập từ dữ liệu thực tế.
\item \textbf{Cơ chế thí điểm và chấp nhận rủi ro:} Nhà nước cho phép \textbf{thí điểm những vấn đề thực tiễn mới} chưa có quy định pháp luật (chưa có lý luận hoàn chỉnh). Việc doanh nghiệp thử nghiệm công nghệ mới dưới sự giám sát của Nhà nước là cách để ``chứng minh chân lý'' trong thực tiễn, giúp hoàn thiện các khung khổ lý luận và pháp lý về công nghệ trong tương lai.
\item \textbf{Đào tạo nguồn nhân lực:} Việc sắp xếp lại hệ thống tổ chức khoa học, công nghệ phải bảo đảm sự \textbf{gắn kết chặt chẽ giữa nghiên cứu -- ứng dụng -- đào tạo}. Nhân lực IT không chỉ cần trình độ lý luận cao (toán học, vật lý, tin học) mà phải có khả năng triển khai các nhiệm vụ trọng điểm quốc gia, đưa tri thức thành sản phẩm công nghệ có giá trị cạnh tranh quốc tế.
\end{itemize}

%% ============================================================
\section*{Câu 10: Giải thích nhận định: Sự vận động phát triển của xã hội loài người trải qua các hình thái kinh tế -- xã hội là một quá trình lịch sử -- tự nhiên; còn học thuyết hình thái kinh tế -- xã hội là ``hòn đá tảng'' của chủ nghĩa duy vật lịch sử.}

\subsection*{1. Tại sao phát triển xã hội là một ``quá trình lịch sử -- tự nhiên''?}

C. Mác khẳng định sự thay thế các hình thái kinh tế -- xã hội (HT KT-XH) tuân theo các quy luật khách quan, giống như các quy luật trong giới tự nhiên, chứ không phụ thuộc vào ý muốn chủ quan của con người hay các lực lượng siêu nhiên.

\begin{itemize}
\item \textbf{Sự chi phối của các quy luật khách quan:} Sự vận động của xã hội không diễn ra ngẫu nhiên mà bị chi phối bởi các quy luật phổ biến, đặc biệt là \textbf{quy luật quan hệ sản xuất phải phù hợp với trình độ phát triển của lực lượng sản xuất}. Khi lực lượng sản xuất phát triển đến một mức độ nhất định, nó mâu thuẫn với quan hệ sản xuất lỗi thời, đòi hỏi một cuộc cách mạng kinh tế để thiết lập quan hệ sản xuất mới, từ đó kéo theo sự thay đổi của toàn bộ kiến trúc thượng tầng.
\item \textbf{Lực lượng sản xuất là nguồn gốc sâu xa:} Mọi sự biến đổi xã hội suy cho cùng đều bắt nguồn từ sự phát triển của \textbf{lực lượng sản xuất} -- kết quả của hoạt động thực tiễn của con người qua các thế hệ.
\item \textbf{Tính thống nhất và đa dạng:}
  \begin{itemize}
  \item \textbf{Thống nhất:} Xu hướng chung của nhân loại là đi từ thấp đến cao qua các hình thái: Cộng sản nguyên thủy $\rightarrow$ Chiếm hữu nô lệ $\rightarrow$ Phong kiến $\rightarrow$ Tư bản chủ nghĩa $\rightarrow$ Cộng sản chủ nghĩa.
  \item \textbf{Đa dạng:} Tùy thuộc vào điều kiện lịch sử, địa lý, văn hóa và nhân tố chủ quan, mỗi dân tộc có con đường phát triển riêng, có thể phát triển \textbf{``rút ngắn''} hoặc bỏ qua một vài hình thái nhất định để tiến lên hình thái cao hơn, nhưng sự bỏ qua này vẫn phải tuân thủ các điều kiện khách quan.
  \end{itemize}
\end{itemize}

\subsection*{2. Tại sao học thuyết HT KT-XH là ``hòn đá tảng'' của chủ nghĩa duy vật lịch sử?}

Học thuyết này được coi là hạt nhân lý luận, là cơ sở vững chắc nhất của toàn bộ quan niệm duy vật về lịch sử vì những lý do sau:

\begin{itemize}
\item \textbf{Chấm dứt sự tùy tiện trong nghiên cứu xã hội:} Trước Mác, các quan niệm về lịch sử thường lộn xộn và duy tâm (coi ý chí của vĩ nhân hay thần linh quyết định). C. Mác đã ``quy những quan hệ xã hội vào những quan hệ sản xuất'' và ``quy những quan hệ sản xuất vào trình độ của những lực lượng sản xuất'', tạo ra cơ sở khoa học để nghiên cứu xã hội chính xác như một ngành khoa học tự nhiên.
\item \textbf{Phân tích xã hội như một chỉnh thể thống nhất:} Học thuyết chỉ ra cấu trúc phức tạp của xã hội bao gồm ba yếu tố cơ bản gắn bó hữu cơ: \textbf{Lực lượng sản xuất} (nền tảng vật chất), \textbf{Quan hệ sản xuất -- Cơ sở hạ tầng} (cơ cấu kinh tế) và \textbf{Kiến trúc thượng tầng} (chính trị, pháp luật, ý thức). Điều này giúp các nhà khoa học không nhìn nhận các sự kiện xã hội một cách rời rạc mà đặt chúng trong mối liên hệ hệ thống.
\item \textbf{Cung cấp phương pháp luận khoa học:} Nó mang lại cho các ngành khoa học xã hội một \textbf{thế giới quan duy vật} và \textbf{phương pháp luận biện chứng} để giải thích các hiện tượng phức tạp như đấu tranh giai cấp, cách mạng xã hội hay sự hình thành các hình thái nhà nước.
\item \textbf{Cơ sở cho chiến lược cách mạng:} Đối với các Đảng Cộng sản, học thuyết này là kim chỉ nam để phân tích thời đại, hoạch định đường lối cải tạo xã hội cũ và xây dựng xã hội mới (chủ nghĩa xã hội) dựa trên sự phát triển của lực lượng sản xuất hiện đại và quan hệ sản xuất tiến bộ.
\item \textbf{Dự báo khoa học về tương lai:} Trên cơ sở phân tích mâu thuẫn nội tại của chủ nghĩa tư bản, học thuyết khẳng định sự diệt vong tất yếu của nó và sự ra đời của xã hội cộng sản như một tất yếu lịch sử.
\end{itemize}

\textbf{Tóm lại}, học thuyết HT KT-XH đã đưa lịch sử trở thành một môn khoa học thực sự bằng cách chỉ ra rằng xã hội không phải là khối kết tinh vững chắc mà là một \textbf{cơ thể sống đang biến đổi liên tục} theo những quy luật tự thân.

%% ============================================================
\section*{Câu 11: Dựa trên điều gì mà Ph.Ăngghen đưa ra nhận định: ``Giống như Charles Darwin đã tìm ra quy luật phát triển của thế giới hữu cơ, Karx Marx đã tìm ra quy luật phát triển của xã hội loài người''.}

Ph.Ăngghen đưa ra nhận định so sánh \textbf{C.Mác} với \textbf{Charles Darwin} dựa trên việc khẳng định C.Mác đã thực hiện một cuộc cách mạng trong khoa học xã hội, tương tự như những gì Darwin đã làm đối với khoa học tự nhiên. Cụ thể, nhận định này dựa trên các cơ sở sau:

\begin{itemize}
\item \textbf{Phát hiện ra quy luật phát triển khách quan của xã hội:} Giống như Darwin đã tìm ra quy luật tiến hóa của các loài trong thế giới hữu cơ, C.Mác đã phát hiện ra quy luật cơ bản chi phối sự vận động và phát triển của lịch sử loài người. C.Mác đã chỉ ra rằng sự phát triển của xã hội không phải là ngẫu nhiên hay tùy tiện mà tuân theo những quy luật khách quan.
\item \textbf{Sáng tạo ra Chủ nghĩa duy vật lịch sử:} Đây được coi là ``thành tựu vĩ đại nhất của tư tưởng khoa học''. Trước C.Mác, việc nghiên cứu lịch sử và chính trị thường rơi vào tình trạng lộn xộn, tùy tiện và duy tâm. C.Mác đã chấm dứt thời kỳ này bằng cách chỉ ra rằng sự phát triển của các \textbf{hình thái kinh tế -- xã hội} là một \textbf{quá trình lịch sử -- tự nhiên}.
\item \textbf{Xác định cơ sở hạ tầng của lịch sử:} C.Mác phát hiện ra tiền đề đầu tiên của lịch sử nhân loại là những con người hiện thực sản xuất ra đời sống vật chất của chính mình. Ông đã ``quy những quan hệ xã hội vào những quan hệ sản xuất, và đem quy những quan hệ sản xuất vào trình độ của những lực lượng sản xuất''.
\item \textbf{Chỉ ra động lực của sự phát triển:} Ăngghen đánh giá cao việc Mác chứng minh được rằng khi \textbf{lực lượng sản xuất} lớn lên, nó sẽ dẫn đến sự thay đổi của \textbf{quan hệ sản xuất}, từ đó làm nảy sinh và phát triển một hình thái tổ chức đời sống xã hội mới cao hơn (ví dụ: từ phong kiến lên tư bản chủ nghĩa).
\item \textbf{Tính khoa học và chính xác:} Nhờ phát kiến này, lần đầu tiên con người có thể nghiên cứu các điều kiện xã hội và sự biến đổi của chúng một cách chính xác như nghiên cứu các khoa học lịch sử tự nhiên.
\end{itemize}

Tóm lại, Ph.Ăngghen đưa ra nhận định trên vì ông coi \textbf{Chủ nghĩa duy vật lịch sử} là một phát minh vạch thời đại, đưa việc nghiên cứu xã hội loài người lên tầm mức khoa học thực sự, giống như cách thuyết tiến hóa đã đưa sinh học lên tầm khoa học.

%% ============================================================
\section*{Câu 12: Chứng minh luận điểm của Ph.Ăngghen: \textit{``Một dân tộc muốn đứng vững trên đỉnh cao của khoa học thì không thể không có tư duy lý luận''... ``Cứ mỗi lần khoa học đạt được thành tựu mới thì triết học phải thay đổi hình thức tồn tại của chính mình''}}

Luận điểm của Ph.Ăngghen về mối quan hệ giữa tư duy lý luận, triết học và khoa học phản ánh sâu sắc tính biện chứng giữa cái chung và cái riêng, giữa sự khái quát hóa và dữ liệu thực chứng. Dưới đây là phần chứng minh hai vế của luận điểm này dựa trên các nguồn tài liệu:

\subsection*{1. ``Một dân tộc muốn đứng vững trên đỉnh cao của khoa học thì không thể không có tư duy lý luận''}

Luận điểm này khẳng định vai trò tiên quyết của tư duy lý luận (mà hạt nhân là triết học) đối với sự phát triển của các ngành khoa học cụ thể vì những lý do sau:

\begin{itemize}
\item \textbf{Cung cấp công cụ tư duy thiết yếu:} Các nhà khoa học tự nhiên không thể tiến lên dù chỉ một bước nếu thiếu các \textbf{phạm trù lôgíc}. Nếu họ phỉ báng hoặc không để ý đến triết học, họ sẽ vô tình trở thành ``nô lệ của những tàn tích thông tục hóa, tồi tệ nhất'' từ các hệ thống triết học lỗi thời.
\item \textbf{Vượt qua khủng hoảng lý luận:} Lịch sử cho thấy khi khoa học đạt đến ngưỡng phát triển mới nhưng bị kìm hãm bởi tư duy cũ (như tư duy siêu hình), nó sẽ rơi vào khủng hoảng. Ph.Ăngghen và V.I. Lênin đều khẳng định \textbf{phép biện chứng} là công cụ duy nhất giúp khoa học tự nhiên vượt qua những khó khăn về mặt lý luận để giải thích đúng đắn các hiện tượng mới.
\item \textbf{Định hướng nghiên cứu:} Triết học đóng vai trò là \textbf{thế giới quan và phương pháp luận} chung nhất, giúp các nhà khoa học xác định hướng nghiên cứu và xây dựng các lý thuyết khoa học hiện đại.
\item \textbf{Tổng hợp tri thức:} Tư duy lý luận giúp kết nối các ngành khoa học chuyên ngành vốn đang phân tán manh mún, từ đó xây dựng nên \textbf{bức tranh khoa học chung về thế giới}.
\end{itemize}

\subsection*{2. ``Cứ mỗi lần khoa học đạt được thành tựu mới thì triết học phải thay đổi hình thức tồn tại của chính mình''}

Vế này phản ánh tính mở và sự phụ thuộc của triết học (đặc biệt là chủ nghĩa duy vật) vào các thành tựu thực chứng của khoa học:

\begin{itemize}
\item \textbf{Khoa học là tiền đề cho sự biến đổi của triết học:} Mỗi phát minh mang ý nghĩa thời đại trong khoa học tự nhiên đều buộc chủ nghĩa duy vật phải \textbf{thay đổi hình thức} để phù hợp với tri thức mới. Triết học không đứng ngoài hay đứng trên khoa học mà phải khái quát hóa những thành tựu mới nhất để hoàn thiện chính mình.
\item \textbf{Minh chứng lịch sử:}
  \begin{itemize}
  \item \textbf{Thuyết nhật tâm của Cô-péc-ních} đã giáng đòn vào thần học, làm hồi sinh và thay đổi diện mạo triết học duy vật thời Phục hưng.
  \item \textbf{Ba phát minh vĩ đại thế kỷ XIX} (Thuyết tế bào, Thuyết tiến hóa, Định luật bảo toàn năng lượng) là cơ sở trực tiếp làm ``vỡ tung'' phương pháp siêu hình, dẫn đến sự ra đời của \textbf{triết học Mác (chủ nghĩa duy vật biện chứng)}.
  \item \textbf{Thuyết tương đối và Cơ học lượng tử} thế kỷ XX đã buộc triết học phải từ bỏ các quan niệm cơ học máy móc về không gian, thời gian và vật chất để tiến tới những khái niệm ``mềm dẻo'', biện chứng hơn.
  \end{itemize}
\item \textbf{Tính tất yếu của sự phát triển:} Triết học Mác là một \textbf{hệ thống lý luận mở}, nó đòi hỏi phải liên tục được bổ sung và phát triển dựa trên thực tiễn sinh động và các kết luận mới từ khoa học để không trở nên lạc hậu, giáo điều.
\end{itemize}

\textbf{Tóm lại:} Mối quan hệ giữa triết học và khoa học là \textbf{quan hệ hai chiều}. Triết học cung cấp ``kim chỉ nam'' về tư duy lý luận để khoa học đạt đỉnh cao; ngược lại, khoa học cung cấp ``nguyên liệu'' thực chứng bắt buộc triết học phải không ngừng đổi mới hình thức tồn tại để phản ánh chân thực thế giới khách quan.

\end{multicols*}
\end{document}
